\appendix

\chapter{Cat\'alogo de fun\c{c}\~oes}
\label{ap:funcoes}

Esta se\c{c}\~ao lista as 129 fun\c{c}\~oes do c\'odigo Rust, agrupadas por
m\'odulo. O \emph{stub} \code{\_start} n\~ao \'e uma fun\c{c}\~ao Rust e
aparece separadamente no Cap\'itulo~\ref{cap:boot}. Declara\c{c}\~oes
\code{extern} sem corpo, como \code{gdt\_reload\_cs}, \code{user\_enter} e
\code{user\_resume}, n\~ao s\~ao contadas.

\section{\code{main.rs} (2 fun\c{c}\~oes)}

\begin{longtable}{@{}p{6.2cm}p{8.3cm}@{}}
\caption{Fun\c{c}\~oes de \code{main.rs}.}\\
\toprule
\textbf{Assinatura} & \textbf{Descri\c{c}\~ao} \\
\midrule
\endfirsthead
\toprule
\textbf{Assinatura} & \textbf{Descri\c{c}\~ao} \\
\midrule
\endhead
\code{kernel\_main(multiboot\_magic: u32, multiboot\_info: u32) -> !} &
Ponto de entrada Rust; inicializa drivers, imprime banner e entra no shell. \\
\code{panic(info: \&PanicInfo) -> !} &
\emph{Panic handler}; imprime a mensagem e para a CPU. \\
\bottomrule
\end{longtable}

\section{\code{arch.rs} (5 fun\c{c}\~oes)}

\begin{longtable}{@{}p{5.2cm}p{9.3cm}@{}}
\caption{Fun\c{c}\~oes de \code{arch.rs}.}\\
\toprule
\textbf{Assinatura} & \textbf{Descri\c{c}\~ao} \\
\midrule
\endfirsthead
\toprule
\textbf{Assinatura} & \textbf{Descri\c{c}\~ao} \\
\midrule
\endhead
\code{inb(port: u16) -> u8} &
L\^e um byte de uma porta. \code{unsafe}. \\
\code{outb(port: u16, value: u8)} &
Escreve um byte em uma porta. \code{unsafe}. \\
\code{enable\_interrupts()} &
Executa \code{sti} para habilitar interrup\c{c}\~oes mascar\'aveis. \\
\code{halt\_loop() -> !} &
La\c{c}o infinito com \code{cli; hlt}. \\
\code{reboot() -> !} &
Envia \code{0xFE} ao 8042; se falhar, para a CPU. \\
\bottomrule
\end{longtable}

\section{\code{sync.rs} (5 fun\c{c}\~oes)}

\begin{longtable}{@{}p{5.7cm}p{8.8cm}@{}}
\caption{Fun\c{c}\~oes de \code{sync.rs}.}\\
\toprule
\textbf{Assinatura} & \textbf{Descri\c{c}\~ao} \\
\midrule
\endfirsthead
\toprule
\textbf{Assinatura} & \textbf{Descri\c{c}\~ao} \\
\midrule
\endhead
\code{SpinLock::new(data: T) -> SpinLock<T>} &
Cria o lock em contexto \code{const}. \\
\code{SpinLock::lock(\&self) -> SpinLockGuard<'\_, T>} &
Adquire o lock girando at\'e conseguir. \\
\code{SpinLockGuard::deref(\&self) -> \&T} &
Acesso imut\'avel ao dado protegido. \\
\code{SpinLockGuard::deref\_mut(\&mut self) -> \&mut T} &
Acesso mut\'avel ao dado protegido. \\
\code{SpinLockGuard::drop(\&mut self)} &
Libera o lock com ordena\c{c}\~ao \code{Release}. \\
\bottomrule
\end{longtable}

\section{\code{vga.rs} (5 fun\c{c}\~oes)}

\begin{longtable}{@{}p{6.1cm}p{8.4cm}@{}}
\caption{Fun\c{c}\~oes de \code{vga.rs}.}\\
\toprule
\textbf{Assinatura} & \textbf{Descri\c{c}\~ao} \\
\midrule
\endfirsthead
\toprule
\textbf{Assinatura} & \textbf{Descri\c{c}\~ao} \\
\midrule
\endhead
\code{blank\_cell() -> u16} &
C\'elula de espa\c{c}o com atributo padr\~ao; \code{const fn}. \\
\code{init()} &
Limpa o hardware e configura a forma do cursor. \\
\code{clear\_hardware()} &
Preenche as 2.000 c\'elulas do framebuffer com c\'elulas vazias. \\
\code{render(cells, row, column)} &
Copia um buffer de console para \code{0xB8000} e posiciona o cursor. \\
\code{set\_cursor(row, column)} &
Escreve a posi\c{c}\~ao nos registradores da CRTC. \\
\bottomrule
\end{longtable}

\section{\code{tty.rs} (31 fun\c{c}\~oes)}

\begin{longtable}{@{}p{6.0cm}p{8.5cm}@{}}
\caption{Fun\c{c}\~oes de \code{tty.rs}.}\\
\toprule
\textbf{Assinatura} & \textbf{Descri\c{c}\~ao} \\
\midrule
\endfirsthead
\toprule
\textbf{Assinatura} & \textbf{Descri\c{c}\~ao} \\
\midrule
\endhead
\code{Tty::new() -> Tty} &
Inicializa um console vazio; \code{const fn}. \\
\code{Tty::clear(\&mut self)} &
Preenche as c\'elulas com espa\c{c}os e zera o cursor. \\
\code{Tty::push\_input(\&mut self, byte: u8)} &
Anexa um byte \`a fila de entrada do console. \\
\code{Tty::pop\_input(\&mut self) -> Option<u8>} &
Remove um byte da fila. \\
\code{Tty::write\_bytes(\&mut self, bytes: \&[u8])} &
Escreve uma fatia no buffer do console. \\
\code{Tty::write\_byte(\&mut self, byte: u8)} &
Despacha texto simples e sequ\^encias de escape ANSI. \\
\code{Tty::write\_plain(\&mut self, byte: u8)} &
Interpreta controles, backspace e ASCII imprim\'ivel. \\
\code{Tty::write\_csi(\&mut self, byte: u8)} &
Acumula par\^ametros CSI e chama o despachante final. \\
\code{Tty::push\_csi\_param(\&mut self)} &
Guarda um par\^ametro decimal da sequ\^encia CSI. \\
\code{Tty::csi\_param(\&self, index, default) -> u16} &
Devolve um par\^ametro ou seu valor padr\~ao. \\
\code{Tty::csi\_count\_param(\&self, index) -> usize} &
Devolve um par\^ametro de contagem, tratando zero como um. \\
\code{Tty::execute\_csi(\&mut self, final\_byte: u8)} &
Executa CUP, movimento do cursor e comandos de apagar. \\
\code{Tty::erase\_display(\&mut self, mode: u16)} &
Implementa ED (\code{ESC[J}, \code{ESC[2J}). \\
\code{Tty::erase\_line(\&mut self, mode: u16)} &
Implementa EL (\code{ESC[K}). \\
\code{Tty::blank\_row(\&mut self, row: usize)} &
Preenche uma linha com c\'elulas vazias. \\
\code{Tty::write\_char(\&mut self, ch: char)} &
Grava uma c\'elula e avan\c{c}a a coluna. \\
\code{Tty::newline(\&mut self)} &
Avan\c{c}a a linha, com rolagem na \'ultima. \\
\code{Tty::backspace(\&mut self)} &
Recua uma posi\c{c}\~ao e apaga a c\'elula. \\
\code{Tty::scroll(\&mut self)} &
Move o conte\'udo uma linha para cima. \\
\code{init()} &
Limpa os oito consoles e renderiza o \code{tty0}. \\
\code{active() -> usize} &
Devolve o \'indice do console ativo. \\
\code{switch(index: usize) -> bool} &
Troca o console ativo e o redesenha. \\
\code{clear()} &
Escreve \code{ESC[2J ESC[H} no console ativo. \\
\code{write\_bytes(bytes: \&[u8])} &
Escreve no console ativo e o renderiza. \\
\code{write\_to(index: usize, bytes: \&[u8])} &
Escreve em um console espec\'ifico sem trocar o ativo. \\
\code{read\_byte() -> Option<u8>} &
Retira um byte da fila do console ativo. \\
\code{read\_byte\_from(index: usize) -> Option<u8>} &
Retira um byte da fila de um console espec\'ifico. \\
\code{pump\_input()} &
Drena teclado e serial para o console ativo; trata
\code{Alt+F1}..\code{Alt+F8}. \\
\code{push\_input(byte: u8)} &
Auxiliar interno que enfileira um byte por console. \\
\code{<TtyWriter as fmt::Write>::write\_str} &
Anexa texto formatado ao console ativo e espelha na serial. \\
\code{\_print(args: fmt::Arguments<'\_>)} &
Ponto de entrada das macros \code{print!}/\code{println!}. \\
\bottomrule
\end{longtable}

\section{\code{keyboard.rs} (10 fun\c{c}\~oes)}

\begin{longtable}{@{}p{5.3cm}p{9.2cm}@{}}
\caption{Fun\c{c}\~oes de \code{keyboard.rs}.}\\
\toprule
\textbf{Assinatura} & \textbf{Descri\c{c}\~ao} \\
\midrule
\endfirsthead
\toprule
\textbf{Assinatura} & \textbf{Descri\c{c}\~ao} \\
\midrule
\endhead
\code{build\_keymap() -> [u8; 128]} &
Constr\'oi a tabela de scancodes em tempo de compila\c{c}\~ao. \\
\code{apply\_shift(byte: u8) -> u8} &
Aplica Shift a letras, d\'igitos e pontua\c{c}\~ao. \\
\code{wait\_input\_clear() -> bool} &
Espera o buffer de entrada esvaziar. \code{unsafe}. \\
\code{wait\_output\_full() -> bool} &
Espera um byte na sa\'ida. \code{unsafe}. \\
\code{flush\_output()} &
Descarta at\'e 64 bytes pendentes. \code{unsafe}. \\
\code{send\_command(command: u8)} &
Envia comando ao 8042. \code{unsafe}. \\
\code{send\_data(data: u8)} &
Envia dado ao teclado. \code{unsafe}. \\
\code{init()} &
Configura tradu\c{c}\~ao, IRQ, clock e varredura. \\
\code{read\_key() -> Option<Key>} &
L\^e e traduz um scancode; devolve um caractere ou pedido de console. \\
\code{read\_char() -> Option<u8>} &
Inv\'olucro de compatibilidade que devolve apenas \code{Key::Char}. \\
\bottomrule
\end{longtable}

\section{\code{serial.rs} (5 fun\c{c}\~oes)}

\begin{longtable}{@{}p{4.8cm}p{9.7cm}@{}}
\caption{Fun\c{c}\~oes de \code{serial.rs}.}\\
\toprule
\textbf{Assinatura} & \textbf{Descri\c{c}\~ao} \\
\midrule
\endfirsthead
\toprule
\textbf{Assinatura} & \textbf{Descri\c{c}\~ao} \\
\midrule
\endhead
\code{init()} &
Configura COM1 em 38400 8N1 e detecta presen\c{c}a. \\
\code{read\_byte() -> Option<u8>} &
Leitura n\~ao bloqueante de um byte. \\
\code{write\_byte(byte: u8)} &
Envia um byte, esperando o THR. Interna. \\
\code{write\_bytes(bytes: \&[u8])} &
Envia bytes com convers\~ao de \code{\textbackslash n} para CRLF. \\
\code{write\_str(text: \&str)} &
Inv\'olucro de \code{write\_bytes} para strings. \\
\bottomrule
\end{longtable}

\section{\code{shell.rs} (9 fun\c{c}\~oes)}

\begin{longtable}{@{}p{4.8cm}p{9.7cm}@{}}
\caption{Fun\c{c}\~oes de \code{shell.rs}.}\\
\toprule
\textbf{Assinatura} & \textbf{Descri\c{c}\~ao} \\
\midrule
\endfirsthead
\toprule
\textbf{Assinatura} & \textbf{Descri\c{c}\~ao} \\
\midrule
\endhead
\code{run() -> !} &
La\c{c}o eterno de prompt, leitura e execu\c{c}\~ao. \\
\code{trim(bytes: \&[u8]) -> \&[u8]} &
Remove espa\c{c}os ASCII das extremidades. \\
\code{as\_text(bytes: \&[u8]) -> \&str} &
Converte bytes em \code{\&str} ou texto de erro. \\
\code{execute(line: \&[u8])} &
Compara e executa o comando. \\
\code{parse\_console(arg: \&[u8]) -> Option<usize>} &
Interpreta um \'indice de console de um d\'igito para \code{tty} e \code{getty}. \\
\code{getty\_demo(console: Option<usize>)} &
Executa o login em ring 3 no console escolhido ou ativo. \\
\code{demand\_demo()} &
Escreve na regi\~ao de demand paging do kernel e mostra o valor. \\
\code{user\_demo(protection: bool)} &
Executa o demo em ring 3 e imprime o status ou a falha. \\
\code{memory\_info()} &
Mostra frames, page faults, ticks e flags U/S das p\'aginas. \\
\bottomrule
\end{longtable}

\section{\code{gdt.rs} (7 fun\c{c}\~oes)}

\begin{longtable}{@{}p{6.0cm}p{8.5cm}@{}}
\caption{Fun\c{c}\~oes de \code{gdt.rs}.}\\
\toprule
\textbf{Assinatura} & \textbf{Descri\c{c}\~ao} \\
\midrule
\endfirsthead
\toprule
\textbf{Assinatura} & \textbf{Descri\c{c}\~ao} \\
\midrule
\endhead
\code{stack\_top<const N: usize>(stack: *mut Stack<N>) -> u64} &
Calcula o topo de uma pilha est\'atica. \\
\code{gdt\_reload\_cs()} &
Recarrega \code{CS} ap\'os carregar a nova GDT. \\
\code{user\_kernel\_stack\_top() -> u64} &
Topo da pilha de 16 KiB usada em interrup\c{c}\~oes vindas de ring 3. \\
\code{double\_fault\_stack\_top() -> u64} &
Topo da pilha IST1 de 8 KiB para \code{\#DF}. \\
\code{descriptor(base, limit, access, flags) -> u64} &
Monta um descritor de segmento de 64 bits. \\
\code{tss\_descriptor(base: u64, limit: u32) -> (u64, u64)} &
Monta as duas entradas de 8 bytes do descritor de TSS. \\
\code{init()} &
Carrega a GDT definitiva, a TSS e as pilhas de privil\'egio. \\
\bottomrule
\end{longtable}

\section{\code{idt.rs} (14 fun\c{c}\~oes)}

\begin{longtable}{@{}p{6.2cm}p{8.3cm}@{}}
\caption{Fun\c{c}\~oes de \code{idt.rs}.}\\
\toprule
\textbf{Assinatura} & \textbf{Descri\c{c}\~ao} \\
\midrule
\endfirsthead
\toprule
\textbf{Assinatura} & \textbf{Descri\c{c}\~ao} \\
\midrule
\endhead
\code{init()} &
Instala as 256 portas, remapeia o PIC e liga o PIT em IRQ0. \\
\code{set\_gate(idt, vector, handler, dpl, ist)} &
Preenche uma porta com seletor de c\'odigo, DPL e \'indice IST. \\
\code{io\_wait()} &
Pequeno atraso de I/O para o PIC assimilar comandos. \\
\code{remap\_pic()} &
Move as IRQs do PIC mestre/escravo para \code{0x20..0x2F}. \\
\code{pit\_init()} &
Programa o divisor do PIT para aproximadamente 100 Hz. \\
\code{eoi(irq: u8)} &
Envia \emph{End of Interrupt} ao PIC mestre e, se preciso, ao escravo. \\
\code{ticks() -> u64} &
Devolve o contador de interrup\c{c}\~oes do timer. \\
\code{exception\_name(vector: u64) -> \&'static str} &
Mapeia o vetor para um nome leg\'ivel de exce\c{c}\~ao. \\
\code{interrupt\_dispatch(frame: *mut InterruptFrame)} &
Ponto \'unico chamado pelos \emph{stubs}; despacha exce\c{c}\~oes, IRQs e a
\emph{syscall}. \\
\code{page\_fault(frame: \&mut InterruptFrame)} &
L\^e \code{CR2} e delega ao \code{paging::handle\_page\_fault}; pode encerrar
o usu\'ario. \\
\code{general\_protection(frame: \&mut InterruptFrame)} &
Trata \code{\#GP}; encerra o usu\'ario ou para o kernel. \\
\code{user\_fault(frame: \&mut InterruptFrame) -> bool} &
Verifica se a falha veio de ring 3 com um programa ativo. \\
\code{syscall(frame: \&mut InterruptFrame)} &
Encaminha \code{int \$0x80} para \code{syscall::dispatch}. \\
\code{fatal(frame: \&InterruptFrame) -> !} &
Imprime o diagn\'ostico da falha e para a CPU. \\
\bottomrule
\end{longtable}

\section{\code{paging.rs} (23 fun\c{c}\~oes)}

\begin{longtable}{@{}p{6.2cm}p{8.3cm}@{}}
\caption{Fun\c{c}\~oes de \code{paging.rs}.}\\
\toprule
\textbf{Assinatura} & \textbf{Descri\c{c}\~ao} \\
\midrule
\endfirsthead
\toprule
\textbf{Assinatura} & \textbf{Descri\c{c}\~ao} \\
\midrule
\endhead
\code{bitmap() -> *mut u8} &
Devolve o ponteiro para o bitmap est\'atico de 32 KiB. \code{unsafe}. \\
\code{frame\_is\_used(frame: usize) -> bool} &
Consulta um bit do bitmap. \\
\code{mark\_used(frame: usize)} &
Marca um frame como em uso. \\
\code{mark\_free(frame: usize)} &
Marca um frame como livre. \\
\code{free\_range(start, end, reserve\_end)} &
Libera uma faixa de RAM dispon\'ivel, preservando o fim do kernel. \\
\code{reserve\_range(start: u64, end: u64)} &
Reserva uma faixa (primeiro MiB, kernel, estruturas Multiboot2). \\
\code{init(multiboot\_magic, multiboot\_info)} &
L\^e o mapa Multiboot2 e inicializa o alocador de frames. \\
\code{alloc\_frame() -> Option<u64>} &
Aloca e zera um frame de 4 KiB; devolve \code{None} se a RAM acabou. \\
\code{free\_frame(frame: u64)} &
Devolve um frame ao bitmap. \\
\code{stats() -> Stats} &
Devolve frames alocados, livres e total. \\
\code{demand\_faults() -> u64} &
Devolve o n\'umero de page faults atendidos por demand paging. \\
\code{read\_cr3() -> u64} &
L\^e o PML4 atual de \code{CR3}. \\
\code{invlpg(virt: u64)} &
Invalida uma entrada da TLB. \code{unsafe}. \\
\code{next\_table(entry, user) -> Result<*mut PageTable, \&'static str>} &
Obt\'em ou cria a pr\'oxima tabela, aplicando U/S quando necess\'ario. \\
\code{map\_page(virt, phys, flags, user\_tables)} &
Mapeia uma p\'agina de 4 KiB e invalida a TLB. \code{unsafe}. \\
\code{page\_flags(virt: u64) -> Option<u64>} &
Devolve a entrada final que traduz um endere\c{c}o, sem criar tabelas. \\
\code{unmap\_page(virt: u64) -> Result<(), \&'static str>} &
Remove um mapeamento e libera o frame f\'isico. \\
\code{ensure\_user\_range(addr, len) -> Result<(), \&'static str>} &
Valida um intervalo do usu\'ario e materializa p\'aginas ausentes com U/S = 1. \\
\code{reset\_user\_demand()} &
Desfaz as p\'aginas de pilha/dados do demo de usu\'ario. \\
\code{paging\_page(address: u64) -> u64} &
Alinha um endere\c{c}o para o in\'icio da p\'agina. \\
\code{handle\_page\_fault(addr, error) -> Result<(), \&'static str>} &
Decide entre materializar uma p\'agina ou recusar o acesso. \\
\code{kernel\_demand\_addr(index: usize) -> u64} &
Endere\c{c}o de uma p\'agina da regi\~ao de demand paging do kernel. \\
\code{install\_user\_code(bytes: \&[u8]) -> Result<(), \&'static str>} &
Copia o programa para \code{USER\_CODE}, criando a p\'agina residente. \\
\bottomrule
\end{longtable}

\section{\code{user.rs} (10 fun\c{c}\~oes)}

\begin{longtable}{@{}p{6.0cm}p{8.5cm}@{}}
\caption{Fun\c{c}\~oes de \code{user.rs}.}\\
\toprule
\textbf{Assinatura} & \textbf{Descri\c{c}\~ao} \\
\midrule
\endfirsthead
\toprule
\textbf{Assinatura} & \textbf{Descri\c{c}\~ao} \\
\midrule
\endhead
\code{Program::new() -> Program} &
Cria um buffer vazio de 64 bytes para o c\'odigo de usu\'ario. \\
\code{Program::push(\&mut self, bytes: \&[u8])} &
Anexa bytes ao programa. \\
\code{Program::movabs\_rax(\&mut self, value: u64)} &
Emite \code{movabs rax, imm64}. \\
\code{program(demo: Demo) -> Program} &
Monta o c\'odigo de m\'aquina dos demos \code{Demand} e \code{Protection}. \\
\code{is\_active() -> bool} &
Indica se h\'a um programa de usu\'ario em execu\c{c}\~ao. \\
\code{run(demo: Demo) -> Result<i32, \&'static str>} &
Instala o c\'odigo, entra em ring 3 e devolve o status de sa\'ida. \\
\code{run\_getty(console: Option<usize>) -> Result<i32, \&'static str>} &
Troca para o console pedido, executa o blob de login e devolve o status. \\
\code{abort(status: i32) -> !} &
Chamado por exce\c{c}\~oes/\emph{syscall}; grava o status e retoma o kernel. \\
\bottomrule
\end{longtable}

\section{\code{syscall.rs} (3 fun\c{c}\~oes)}

\begin{longtable}{@{}p{5.6cm}p{8.9cm}@{}}
\caption{Fun\c{c}\~oes de \code{syscall.rs}.}\\
\toprule
\textbf{Assinatura} & \textbf{Descri\c{c}\~ao} \\
\midrule
\endfirsthead
\toprule
\textbf{Assinatura} & \textbf{Descri\c{c}\~ao} \\
\midrule
\endhead
\code{dispatch(frame: \&mut InterruptFrame)} &
Despacha \code{exit}, \code{write}, \code{read} e \code{get\_tty}. \\
\code{sys\_write(frame: \&mut InterruptFrame)} &
Valida o buffer do usu\'ario, escreve no console ativo e devolve o tamanho. \\
\code{sys\_read(frame: \&mut InterruptFrame)} &
Valida o buffer e l\^e uma linha do console, com eco opcional. \\
\bottomrule
\end{longtable}

\chapter{Mapa de portas e mem\'oria}
\label{ap:portas}

\section{Portas de I/O}

\begin{table}[H]
\centering
\caption{Portas de I/O usadas pelo kernel.}
\label{tab:portmap}
\begin{tabular}{@{}llp{7.5cm}@{}}
\toprule
\textbf{Porta} & \textbf{Dispositivo} & \textbf{Uso} \\
\midrule
\code{0x60} & 8042 dados & Scancodes e dados ao teclado \\
\code{0x64} & 8042 status/comando & Status, comandos, reset \\
\code{0x3D4} & VGA CRTC \'indice & Sele\c{c}\~ao de registrador \\
\code{0x3D5} & VGA CRTC dados & Posi\c{c}\~ao e forma do cursor \\
\code{0x20--0x21} & PIC mestre & Comandos, m\'ascara e EOI \\
\code{0xA0--0xA1} & PIC escravo & Comandos, m\'ascara e EOI \\
\code{0x40--0x43} & PIT 8253/8254 & Canal 0, modo e divisor do timer \\
\code{0x3F8--0x3FF} & UART 16550 COM1 & Dados, LSR, divisor, FIFO \\
\bottomrule
\end{tabular}
\end{table}

\section{Mapa de mem\'oria}

\begin{table}[H]
\centering
\caption{Regi\~oes de mem\'oria relevantes.}
\label{tab:memmap}
\begin{tabularx}{\linewidth}{@{}lX@{}}
\toprule
\textbf{Endere\c{c}o} & \textbf{Uso} \\
\midrule
\code{0x00000000--0x000FFFFF} & Mem\'oria baixa, BIOS, estruturas do GRUB. \\
\code{0x00100000} & In\'icio de \code{.boot} (stub, cabe\c{c}alho, GDT,
tabelas de p\'agina, pilha de 8 KiB). \\
\code{0x00107000} & \emph{Load address} de \code{.text} (virtual alto). \\
\code{0x00109518} & \emph{Load address} de \code{.rodata}. \\
\code{0x0010A000} & \emph{Load address} de \code{.data}. \\
\code{0x0010A020} & \emph{Load address} de \code{.bss}; inclui a pilha de
32 KiB do kernel. \\
\code{0x000B8000} & Framebuffer de texto VGA (80x25, 4.000 bytes). \\
\code{0xFFFFFFFF80000000} & Base virtual do kernel (metade alta). \\
\code{0xFFFFFFFF80107000} & In\'icio virtual de \code{.text}. \\
\code{0x0000008000001000} & \code{USER\_CODE}: p\'agina residente do demo. \\
\code{0x0000008000002000} & \code{USER\_DATA}: p\'agina sob demanda. \\
\code{0x0000008000007FF8} & Topo da pilha do usu\'ario; a p\'agina abaixo
\'e sob demanda. \\
\code{0xFFFFFFFFC0000000} & \code{KDEMAND\_BASE}: 2 MiB de demand paging do
kernel. \\
\bottomrule
\end{tabularx}
\end{table}

\begin{nota}[Endere\c{c}os observados]
Os valores de \code{.text}, \code{.rodata}, \code{.data} e \code{.bss}
correspondem \`a compila\c{c}\~ao de refer\^encia. Eles podem mudar com
otimiza\c{c}\~oes, vers\~ao do compilador ou altera\c{c}\~oes no c\'odigo.
O que n\~ao muda \'e o esquema: stub em 1 MiB, kernel em
\code{0xFFFFFFFF80000000}.
\end{nota}

\chapter{Gloss\'ario}
\label{ap:glossario}

\begin{description}[style=nextline,leftmargin=2em]
  \item[BIOS] Firmware legado de inicializa\c{c}\~ao de PCs, respons\'avel
        por detectar hardware e iniciar o bootloader.
  \item[Bootloader] Programa que carrega o kernel; aqui, o GRUB.
  \item[CR0/CR3/CR4] Registradores de controle da CPU x86; controlam
        pagina\c{c}\~ao, modo protegido e features.
  \item[CRTC] Controlador de tubo de raios cat\'odicos da VGA; controla
        cursor e timing de v\'ideo.
  \item[El Torito] Padr\~ao para tornar CDs/DVDs inicializ\'aveis.
  \item[Demand paging] T\'ecnica de criar a p\'agina somente no primeiro
        acesso, quando a CPU gera \code{\#PF}.
  \item[GDT] Global Descriptor Table; define segmentos de c\'odigo e dados.
  \item[GOP] Graphics Output Protocol, interface de framebuffer em UEFI.
  \item[IDT] Interrupt Descriptor Table; define handlers de interrup\c{c}\~oes
        e exce\c{c}\~oes.
  \item[Page fault] Exce\c{c}\~ao (\code{\#PF}) gerada quando uma p\'agina
        n\~ao est\'a presente ou viola as permiss\~oes U/S.
  \item[PIC] Programmable Interrupt Controller 8259; prioriza e entrega as
        IRQs \`a CPU.
  \item[PIT] Programmable Interval Timer 8253/8254; gera a IRQ0 do timer.
  \item[Ring 3] N\'ivel de privil\'egio de usu\'ario; c\'odigo com acesso
        restrito \`a mem\'oria marcada U/S.
  \item[TSS] Task State Segment; fornece \code{rsp0} e pilhas IST ao trocar
        de ring 3 para ring 0.
  \item[LME] Long Mode Enable, bit do registrador EFER que habilita o modo
        de 64 bits.
  \item[LSR] Line Status Register da UART; indica dados dispon\'iveis e
        transmissor vazio.
  \item[Multiboot2] Especifica\c{c}\~ao de protocolo entre bootloader e
        kernel; usada pelo GRUB.
  \item[PML4/PDPT/PD] N\'iveis da tabela de p\'aginas x86\_64.
  \item[Polling] Consulta repetida de hardware, sem interrup\c{c}\~oes.
  \item[Ring 0] N\'ivel de privil\'egio mais alto da CPU; c\'odigo de kernel.
  \item[Scancode] C\'odigo enviado pelo teclado para representar uma tecla.
  \item[Spinlock] Lock que faz a CPU girar at\'e a regi\~ao cr\'itica ser
        liberada.
  \item[Stub] C\'odigo inicial em assembly que prepara a CPU antes do Rust.
  \item[Triple fault] Falha de exce\c{c}\~ao durante o tratamento de outra,
        que faz a CPU resetar.
  \item[UART] Universal Asynchronous Receiver/Transmitter; controlador
        serial.
  \item[VGA texto] Modo de v\'ideo legado com matriz de caracteres em
        \code{0xB8000}.
\end{description}

\begin{thebibliography}{9}

\bibitem{intel}
Intel Corporation.
\newblock \emph{Intel 64 and IA-32 Architectures Software Developer's
Manual}.
\newblock Volume 3A: System Programming Guide, cap\'itulos sobre modo
protegido, pagina\c{c}\~ao e long mode.

\bibitem{multiboot2}
Free Software Foundation.
\newblock \emph{Multiboot2 Specification}.
\newblock Dispon\'ivel em \url{https://www.gnu.org/software/grub/manual/multiboot2/}.

\bibitem{linux}
Linux Kernel Organization.
\newblock \emph{Linux Kernel 7.3.1 Source Tree}.
\newblock Arquivos \code{arch/x86}, \code{drivers/video/console/vgacon.c},
\code{drivers/input}, \code{drivers/tty/serial/8250} e \code{kernel/}.

\bibitem{osdev}
OSDev Wiki.
\newblock \emph{PS/2 Keyboard}, \emph{Serial Ports}, \emph{VGA Text Mode} e
\emph{Setting Up Long Mode}.
\newblock Dispon\'ivel em \url{https://wiki.osdev.org/}.

\bibitem{rust}
The Rust Project.
\newblock \emph{The Rust Programming Language} e \emph{Rust Platform Support:
x86\_64-unknown-none}.
\newblock Dispon\'ivel em \url{https://doc.rust-lang.org/}.

\bibitem{qemu}
QEMU Project.
\newblock \emph{QEMU Emulator User Documentation}.
\newblock Dispon\'ivel em \url{https://www.qemu.org/documentation/}.

\bibitem{grub}
GNU GRUB.
\newblock \emph{GNU GRUB Manual}.
\newblock Dispon\'ivel em \url{https://www.gnu.org/software/grub/manual/grub/}.

\end{thebibliography}
